% =====================================================================
%  CHAPTER 6: 儿童少年心理行为发育（对应大纲第五章：儿童心理发育）
%  参考往届笔记：第四章 儿童少年心理行为发育
% =====================================================================
\chapter{儿童少年心理行为发育}
\setcounter{chapter}{6}
\chapterintro{
\textbf{大纲要求：}
\begin{itemize}[leftmargin=*, itemsep=2pt]
    \item \textbf{掌握}（双横线）：儿童心理发育的特点；儿童少年认知能力发育
    \item \textbf{熟悉}（单横线）：儿童少年情绪情感发育；儿童少年个性及社会化发展
\end{itemize}
\vspace{4pt}
本章介绍儿童少年心理行为发育的全过程。心理行为发育包括认知、语言、情绪情感、人格和社会化适应等多个方面，其发展趋势是从简单到复杂、从低级到高级的上升过程。
}

% ----- 6.0 心理发育概述 -----
\section{儿童心理发育的特点}

\begin{keybox}
\textbf{心理行为发育（Psychological and Behavior Development）：}从受精卵开始到出生、成熟，儿童心理和行为特征的变化过程，包括认知、语言、情绪情感、人格和社会化适应等多个方面。

\textbf{儿童心理发育的特点：}在整个发育过程中，心理行为发育的总趋势是从简单到复杂、从低级到高级的上升过程，各种心理过程和行为特征相互联系、相互影响、共同发展。
\end{keybox}

% ----- 6.1 认知能力发育 -----
\section{儿童少年认知能力发育}

\begin{keybox}
\textbf{认知（Cognition）：}认识活动或认识过程，是大脑反映客观事物的特征、状态及其相互关系、并揭示事物对人的意义与作用的一种高级心理活动。
\end{keybox}

\subsection{感知觉发育}

\begin{keybox}
\textbf{感知觉（Sensory Perception）：}人脑对当前作用于感觉器官的客观事物的反映。

\textbf{一、感觉（Sense）：}对事物个别属性的反映，依赖于个别感觉器官的活动，分为视觉、听觉、嗅觉、味觉和触觉。

\begin{enumerate}[leftmargin=*]
    \item \textbf{视觉}
    \begin{itemize}[leftmargin=*]
        \item 2个月：视觉集中明显形成（尤其是人脸）。
        \item 4个月：出现辨色视觉。
        \item 6个月以前：视力发展敏感期，视力和立体视觉都逐渐发育。
        \item 2$\sim$3岁：能正确辨别红、黄、绿、蓝4种基本颜色并出现双眼视觉。
        \item 出生后3年内：视觉系统功能基本发育成熟。
    \end{itemize}

    \item \textbf{听觉}
    \begin{itemize}[leftmargin=*]
        \item 出生时：听觉器官发育基本成熟，新生儿能听见声音及区分声音的频率、强度和持续时间。
        \item 2岁以后：听力已经很灵敏，几乎达到成人水平。
        \item 学龄初期儿童：听觉感受性随着年龄增长不断发展，在音调辨别、言语听觉、语音听觉和听觉敏感等方面明显提高。
        \item 青少年时期：听觉能力已基本达到稳定水平。
    \end{itemize}

    \item \textbf{嗅觉}
    \begin{itemize}[leftmargin=*]
        \item 新生儿：嗅觉与成人一样敏感，出生时嗅觉发育已比较成熟，能表现出对不同气味的反应。
        \item 7$\sim$8个月：嗅觉逐渐灵敏，能分辨出芳香气味。
        \item 2岁左右：已能很好辨别各种气味。
    \end{itemize}

    \item \textbf{味觉}
    \begin{itemize}[leftmargin=*]
        \item 味觉在儿童时期最发达，以后逐渐衰退。
        \item 2h的新生儿：能分辨甜、酸、苦、咸等多种味道。
        \item 6个月$\sim$1岁：幼儿在这一阶段味觉发展最灵敏。
    \end{itemize}

    \item \textbf{触觉}
    \begin{itemize}[leftmargin=*]
        \item 触觉是人体发展最早、最基本的感觉；触觉在胎儿期已开始发育，到新生儿阶段已高度敏感，触觉是婴儿认识世界的主要手段之一；对婴儿的抚触就是通过对触觉的刺激，增强婴儿触觉的敏感性，加强对外界的反应，促进其发育。
        \item 2$\sim$3岁：能很好辨别各种物体的不同属性。
        \item 5$\sim$6岁：皮肤觉分辨的精细度逐渐提高，在学龄前期已逐步趋于完善。
    \end{itemize}
\end{enumerate}

\textbf{二、知觉（Perception）：}人对事物整体属性的综合反映，是在感觉基础上发展起来的，依赖多种感觉器官，分为时间知觉、空间知觉和运动知觉。

\begin{itemize}[leftmargin=*]
    \item 婴儿3个月时有了分辨简单形状的能力，6个月以前能辨别大小，9$\sim$12个月能知觉形状。
    \item 随着感觉功能的完善促使幼儿的感知觉进一步发展，主要表现为在空间知觉上辨别形状的能力逐渐增强。
    \item 学龄期儿童的视觉感受性和听觉感受性会随年龄增长而不断发展，空间知觉、方位知觉、距离知觉和时间知觉都有很大进步。
\end{itemize}
\end{keybox}

\subsection{注意力和记忆力}

\begin{keybox}
\textbf{一、注意（Attention）：}心理活动对一定对象有选择的集中，是认知过程的开始，分为无意注意和有意注意。

\begin{itemize}[leftmargin=*]
    \item 学龄期儿童：有意注意进一步发展，注意的稳定性随年龄增长而增强，注意集中时间也随之提高。
    \item 青少年期：注意的集中性和稳定性不断增长，稳定性大概能维持40min。
\end{itemize}

\textbf{二、记忆（Memory）：}人脑对过去经验的反映，包括识记、保持和再现（回忆和再认），需经历感知觉记忆、短时记忆和长时记忆3个阶段。

\begin{itemize}[leftmargin=*]
    \item 婴幼儿时期：记忆迅速发展的第一个时期，条件反射的出现是记忆发生的标志。记忆容量增加显著，但记忆以无意识记为主，有意识记成分在逐渐增加。
    \item 学龄期：由于系统学习需要记住大量东西，此时记忆发生质的变化。
    \item 青少年期：记忆的整体水平处于人生的最佳时期。有意识记日益占主导地位，对抽象材料的记忆能力也明显增强。
\end{itemize}
\end{keybox}

\subsection{思维和想象能力}

\begin{keybox}
\textbf{一、思维（Thinking）：}人脑对客观事物间接、概括的反映，能认识事物的本质和事物间的内在联系，属认知的高级阶段。

\begin{itemize}[leftmargin=*]
    \item 婴幼儿期：思维依靠感知觉和动作完成；思维主要是直觉行动思维，是指婴幼儿在动作中进行的思维。婴幼儿在进行这种思维时，只能反映自身动作所触及的具体事物，而不能离开动作，在感知和动作外思考。
    \item 学龄前期：直观行动思维$\to$具体形象思维$\to$抽象逻辑思维，其中占主导地位的是具体形象思维。
    \item 学龄期：思维的基本特点是从具体形象思维逐步过渡到抽象逻辑思维，这是思维发展过程中质的变化。
    \item 青春期：思维能力有很大发展，抽象逻辑思维处于优势地位，能运用抽象思维突破心理运算的界限和理解各种抽象概念，并获得更多增长新知识的机会。
\end{itemize}

\textbf{二、想象（Imagine）：}对已有表象进行加工改造，形成新形象的过程。思维是想象的基础。

\begin{itemize}[leftmargin=*]
    \item 婴幼儿期：1$\sim$2岁的婴幼儿出现想象的萌芽，主要通过动作和口头言语表达出来；3岁时，随着经验与言语的发育，逐渐产生了具有简单想象的游戏。
    \item 学龄前期：已具有丰富的想象力，集中表现在游戏活动中，非游戏时想象水平较低。
    \item 学龄期：想象的有意性、目的性逐渐增强，创造性成分逐渐增大，现实性逐渐提高。
    \item 青少年期：想象能力随教学活动的深入进一步发展，主要表现为有意想象占主要地位，创造想象日益占优势地位，再创想象能力更加独立、概括和精确，想象内容更加复杂，抽象概括性和逻辑性更高。
\end{itemize}
\end{keybox}

\subsection{语言能力发展}

\begin{keybox}
\textbf{语言（Language）：}人类社会中客观存在的现象，是一种社会上约定俗成的符号系统，是由词汇（包括音形义）按一定语法构成的。

\begin{enumerate}[leftmargin=*]
    \item \textbf{前语言阶段（Prespeech Stage）：}从婴儿出生到第一个有真正意义的词产生前的这一时期。
    \item \textbf{语音的发展：}学龄前儿童发音的正确率随年龄的增长而提高，对韵母的发音比较容易掌握，发音错误多出现于辅音中的翘舌音和齿音。
\end{enumerate}
\end{keybox}

% ----- 6.2 情绪情感发育 -----
\section{儿童少年情绪情感发育}

\begin{defbox}
\textbf{一、概念}

\textbf{情绪（Emotion）：}客观事物是否符合人的需要产生的态度体验，反映客观事物与人的需要间的关系。

\textbf{情感（Emotional Feeling）：}与人的高级社会性需要满足与否相联系的态度体验，是在社会交往的实践中逐渐形成的，如友谊感、道德感、美感和理智感等，是人类独有的一种情绪状态。

\textbf{二、儿童青少年情绪情感发育特点}

\begin{enumerate}[leftmargin=*]
    \item \textbf{婴幼儿情绪情感发育特点：}
    \begin{itemize}[leftmargin=*]
        \item 与生理需求是否得到满足直接相关。
        \item 是儿童与生俱来的遗传本能，有先天性。
    \end{itemize}

    \item \textbf{学龄前儿童情绪情感发育特点：}
    \begin{itemize}[leftmargin=*]
        \item 情绪的冲动性逐渐减少。
        \item 情绪情感以外显性为主，内隐性逐渐增强。
        \item 情绪情感以易变性为主，稳定性逐渐发展。
    \end{itemize}

    \item \textbf{学龄期儿童情绪情感发育特点：}
    \begin{itemize}[leftmargin=*]
        \item 情感内容逐渐丰富。
        \item 情感的深刻性和稳定性逐渐增加。
        \item 高级情感逐渐发展。
    \end{itemize}

    \item \textbf{青少年情绪情感发育特点：}
    \begin{itemize}[leftmargin=*]
        \item 情绪情感迅速而热烈，有两极性。
        \item 内隐性和表现性共同存在。
        \item 高级情感已经有了相当的发展。
    \end{itemize}
\end{enumerate}
\end{defbox}

% ----- 6.3 个性及社会化发展 -----
\section{儿童少年个性及社会化发展}

\subsection{个性的发展}

\begin{defbox}
\textbf{一、个性（Personality）/人格：}个人的整个精神面貌，即具有一定倾向性的心理特征的总和，包括一个人个性倾向性、个性心理特征及自我意识。

\textbf{二、气质（Temperament）：}婴幼儿期出生后最早表现出来的一种较为明显而稳定的个性特征，在婴幼儿社会性发展过程中有重要的地位和作用。
\end{defbox}

\subsection{自我意识的发展}

\begin{defbox}
\textbf{自我意识（Self-Consciousness）：}由自我概念、自我评价和自我（情绪）体验等共同组成的心理过程，也是个体不断社会化的过程。

\begin{itemize}[leftmargin=*]
    \item \textbf{自我概念（Self-Concept）：}个人心目中对自己的印象，包括对自己存在、个人身体能力、性格、态度、思想等方面的认识。
    \item \textbf{自我评价（Self-Evaluation）：}自我意识发展的主要成分和标志。
\end{itemize}
\end{defbox}

\subsection{社会化的发展}

\begin{defbox}
\textbf{社会化（Socialization）：}个体在特定社会文化环境中，形成适合于该社会与文化的人格，掌握该社会公认的行为方式，成为合格社会成员的过程。
\end{defbox}

% ----- 复习题 -----
\section{复习题}

\begin{enumerate}[leftmargin=*, itemsep=8pt]
    \item \textbf{【名词解释】}心理行为发育；认知；感知觉；感觉；知觉
    \item \textbf{【名词解释】}注意；记忆；思维；想象；语言
    \item \textbf{【名词解释】}情绪；情感；个性；气质；自我意识；社会化
    \item \textbf{【简答题】}简述儿童心理发育的主要特点。
    \item \textbf{【简答题】}简述儿童少年视觉、听觉、嗅觉、味觉和触觉的发育特点。
    \item \textbf{【简答题】}简述儿童注意力发展的年龄特点。
    \item \textbf{【简答题】}简述儿童记忆发展的阶段特征。
    \item \textbf{【简答题】}简述儿童思维发展的阶段特征。
    \item \textbf{【简答题】}简述儿童想象发展的阶段特征。
    \item \textbf{【简答题】}简述儿童青少年情绪情感发育的阶段性特点。
    \item \textbf{【简答题】}简述个性、气质和自我意识的概念及关系。
    \item \textbf{【简答题】}简述社会化的概念。
\end{enumerate}

\subsection*{参考答案要点}

\begin{enumerate}[leftmargin=*, itemsep=5pt]
    \item \textbf{心理行为发育}：从受精卵开始到出生、成熟，儿童心理和行为特征的变化过程，包括认知、语言、情绪情感、人格和社会化适应等多个方面。\textbf{认知}：认识活动或认识过程，是大脑反映客观事物的特征、状态及其相互关系、并揭示事物对人的意义与作用的一种高级心理活动。\textbf{感知觉}：人脑对当前作用于感觉器官的客观事物的反映。\textbf{感觉}：对事物个别属性的反映，依赖于个别感觉器官的活动，分为视觉、听觉、嗅觉、味觉和触觉。\textbf{知觉}：人对事物整体属性的综合反映，在感觉基础上发展起来，依赖多种感觉器官，分为时间知觉、空间知觉和运动知觉。

    \item \textbf{注意}：心理活动对一定对象有选择的集中，是认知过程的开始，分为无意注意和有意注意。\textbf{记忆}：人脑对过去经验的反映，包括识记、保持和再现（回忆和再认），需经历感知觉记忆、短时记忆和长时记忆3个阶段。\textbf{思维}：人脑对客观事物间接、概括的反映，能认识事物的本质和事物间的内在联系，属认知的高级阶段。\textbf{想象}：对已有表象进行加工改造，形成新形象的过程，思维是想象的基础。\textbf{语言}：人类社会中客观存在的现象，是一种社会上约定俗成的符号系统，由词汇（包括音形义）按一定语法构成。

    \item \textbf{情绪}：客观事物是否符合人的需要产生的态度体验，反映客观事物与人的需要间的关系。\textbf{情感}：与人的高级社会性需要满足与否相联系的态度体验，在社会交往的实践中逐渐形成，如友谊感、道德感、美感和理智感等，是人类独有的一种情绪状态。\textbf{个性/人格}：个人的整个精神面貌，即具有一定倾向性的心理特征的总和，包括一个人个性倾向性、个性心理特征及自我意识。\textbf{气质}：婴幼儿期出生后最早表现出来的一种较为明显而稳定的个性特征，在婴幼儿社会性发展过程中有重要的地位和作用。\textbf{自我意识}：由自我概念、自我评价和自我（情绪）体验等共同组成的心理过程，也是个体不断社会化的过程。\textbf{社会化}：个体在特定社会文化环境中，形成适合于该社会与文化的人格，掌握该社会公认的行为方式，成为合格社会成员的过程。

    \item 儿童心理发育的特点：心理行为发育的总趋势是从简单到复杂、从低级到高级的上升过程，各种心理过程和行为特征相互联系、相互影响、共同发展。包括认知、语言、情绪情感、人格和社会化适应等多个方面。

    \item (1)视觉：2个月视觉集中明显形成（尤其是人脸）；4个月出现辨色视觉；6个月以前为视力发展敏感期；2$\sim$3岁能正确辨别红黄绿蓝4种基本颜色并出现双眼视觉；出生后3年内视觉系统功能基本发育成熟。(2)听觉：出生时听觉器官发育基本成熟，新生儿能区分声音的频率、强度和持续时间；2岁以后听力几乎达到成人水平；学龄初期听觉感受性不断提高；青少年时期听觉能力基本稳定。(3)嗅觉：新生儿嗅觉与成人一样敏感，出生时已比较成熟；7$\sim$8个月能分辨芳香气味；2岁左右能很好辨别各种气味。(4)味觉：儿童时期最发达，以后逐渐衰退；2h新生儿能分辨甜酸苦咸等多种味道；6个月$\sim$1岁味觉发展最灵敏。(5)触觉：人体发展最早、最基本的感觉，胎儿期已开始发育，新生儿阶段已高度敏感，是婴儿认识世界的主要手段之一；2$\sim$3岁能很好辨别各种物体的不同属性；5$\sim$6岁皮肤觉分辨精细度逐渐提高，学龄前期已逐步趋于完善。

    \item 学龄期儿童有意注意进一步发展，注意的稳定性随年龄增长而增强，注意集中时间也随之提高。青少年期注意的集中性和稳定性不断增长，稳定性大概能维持40min。

    \item (1)婴幼儿时期：记忆迅速发展的第一个时期，条件反射的出现是记忆发生的标志，记忆容量增加显著，但记忆以无意识记为主，有意识记成分在逐渐增加。(2)学龄期：由于系统学习需要记住大量东西，此时记忆发生质的变化。(3)青少年期：记忆的整体水平处于人生的最佳时期，有意识记日益占主导地位，对抽象材料的记忆能力也明显增强。

    \item (1)婴幼儿期：思维依靠感知觉和动作完成，主要是直觉行动思维，只能反映自身动作所触及的具体事物，不能离开动作在感知和动作外思考。(2)学龄前期：直观行动思维$\to$具体形象思维$\to$抽象逻辑思维，其中占主导地位的是具体形象思维。(3)学龄期：思维的基本特点是从具体形象思维逐步过渡到抽象逻辑思维，这是思维发展过程中质的变化。(4)青春期：思维能力有很大发展，抽象逻辑思维处于优势地位，能运用抽象思维突破心理运算的界限和理解各种抽象概念。

    \item (1)婴幼儿期：1$\sim$2岁的婴幼儿出现想象的萌芽，主要通过动作和口头言语表达出来；3岁时逐渐产生了具有简单想象的游戏。(2)学龄前期：已具有丰富的想象力，集中表现在游戏活动中，非游戏时想象水平较低。(3)学龄期：想象的有意性、目的性逐渐增强，创造性成分逐渐增大，现实性逐渐提高。(4)青少年期：有意想象占主要地位，创造想象日益占优势地位，再创想象能力更加独立、概括和精确，想象内容更加复杂，抽象概括性和逻辑性更高。

    \item (1)婴幼儿：与生理需求是否得到满足直接相关；是儿童与生俱来的遗传本能，有先天性。(2)学龄前：情绪的冲动性逐渐减少；情绪情感以外显性为主，内隐性逐渐增强；情绪情感以易变性为主，稳定性逐渐发展。(3)学龄期：情感内容逐渐丰富；情感的深刻性和稳定性逐渐增加；高级情感逐渐发展。(4)青少年：情绪情感迅速而热烈，有两极性；内隐性和表现性共同存在；高级情感已经有了相当的发展。

    \item \textbf{个性/人格}：个人的整个精神面貌，即具有一定倾向性的心理特征的总和，包括个性倾向性、个性心理特征及自我意识。\textbf{气质}：婴幼儿期出生后最早表现出来的一种较为明显而稳定的个性特征，在婴幼儿社会性发展过程中有重要的地位和作用。\textbf{自我意识}：由自我概念、自我评价和自我（情绪）体验等共同组成的心理过程，也是个体不断社会化的过程；自我概念是个人心目中对自己的印象，自我评价是自我意识发展的主要成分和标志。三者关系：气质是个性中最早表现出来的稳定特征，自我意识是个性发展的核心。

    \item 社会化是个体在特定社会文化环境中，形成适合于该社会与文化的人格，掌握该社会公认的行为方式，成为合格社会成员的过程。
\end{enumerate}

\newpage
